% Volume three architecture figures. Independent export only; not book fonts.
\colorlet{MutedBlueFill}{FigureBlueFill}
\colorlet{MutedGreenFill}{FigureGreenFill}
\colorlet{MutedPurple}{FigurePurple}
\colorlet{MutedPurpleFill}{FigurePurpleFill}
\tikzset{
 v3 picture/.style={x=1mm,y=1mm,font=\figurefont\fontsize{8}{10}\selectfont,text=FigureInk,line width=.65pt},
 v3 box/.style={draw=FigureBlue,fill=white,rounded corners=1.8pt,line width=.7pt,inner sep=2.5pt,align=center},
 v3 data/.style={v3 box,fill=FigureBlueFill},
 v3 model/.style={v3 box,draw=FigureGreen,fill=FigureGreenFill},
 v3 output/.style={v3 box,draw=FigureYellow,fill=FigureYellowFill},
 v3 risk/.style={v3 box,draw=FigureRed,fill=FigureRedFill},
 v3 group/.style={draw=FigureBlue,rounded corners=1.8pt,line width=.65pt,inner sep=3mm},
 v3 flow/.style={draw=FigureStroke,line width=.75pt,-{Stealth[length=1.65mm,width=1.05mm]},rounded corners=1pt},
 v3 guide/.style={draw=FigureMuted,line width=.6pt,dashed},
 v3 feedback/.style={v3 flow,draw=FigureGreen,dashed},
 v3 title/.style={draw=none,fill=none,font=\figurefont\fontsize{8}{10}\selectfont,align=center,inner sep=1pt},
 v3 note/.style={draw=none,fill=none,font=\FigureNoteFont\fontsize{7.5}{9.5}\selectfont,text=FigureMuted,align=center,inner sep=1pt},
 v3 label/.style={draw=none,fill=none,font=\figurefont\fontsize{8}{10}\selectfont,align=center,inner sep=1pt},
 v3 wire/.style={draw=FigureMuted,line width=.45pt},
 v3 point/.style={circle,draw=FigureBlue,fill=FigureBlueFill,minimum size=2.3mm,inner sep=0pt,line width=.55pt}
}
% Bring pre-existing native TikZ geometry to the same typography and scale.
\tikzset{
 foundation picture/.append style={font=\figurefont\fontsize{8}{10}\selectfont,line width=.65pt},
 foundation node/.append style={rounded corners=1.8pt,line width=.7pt,font=\figurefont\fontsize{8}{10}\selectfont},
 foundation heading/.append style={draw=none,fill=none,font=\figurefont\fontsize{8}{10}\selectfont},
 foundation note/.append style={draw=none,fill=none,font=\FigureNoteFont\fontsize{7.5}{9.5}\selectfont},
 foundation flow/.append style={line width=.75pt,-{Stealth[length=1.65mm,width=1.05mm]}},
 foundation frame/.append style={rounded corners=1.8pt,line width=.7pt},
 foundation edge label/.append style={draw=none,fill=none,font=\figurefont\fontsize{8}{10}\selectfont}
}
% Official Lucide outline icons, consistently used across diagrams. Size in mm.
\newcommand{\VIcon}[5]{\node[inner sep=0pt] at (#2,#3) {\includegraphics[width=#4mm]{figures/03-foundations/00-shared/tikz/icons/#1-#5.pdf}};}
% A neutral network glyph (not a claim about an actual model architecture).
% Coordinates denote its centre; scale 1 gives 13 x 8 mm.
\newcommand{\VNet}[4]{\begin{scope}[shift={(#1,#2)},scale=#3]
 \foreach \a in {-3,0,3}{\foreach \b in {-3,0,3}{\draw[v3 wire] (-5,\a)--(0,\b);}}
 \foreach \a in {-3,0,3}{\draw[v3 wire] (0,\a)--(5,0);}
 \foreach \a in {-3,0,3}{\draw[draw=#4,fill=#4,line width=.55pt] (-5,\a) circle(.9);\draw[draw=#4,fill=white,line width=.55pt] (0,\a) circle(.9);}
 \draw[draw=#4,fill=#4,line width=.55pt] (5,0) circle(.9);
\end{scope}}
\newcommand{\VDots}[4]{\foreach \k in {0,1,2,3}{\fill[#4] ({#1+2.7*\k},#2) circle(#3);}}
% Feature vector/tensor elements are scientific glyphs, not decorative icons.
\newcommand{\VVector}[4]{\foreach \k in {0,...,4}{\fill[#4,rounded corners=.25pt] ({#1+2.2*\k},#2) rectangle ++(1.5,#3);}}

\tikzset{v3 flow/.append style={line width=.95pt,rounded corners=1.8pt},v3 box/.append style={line width=.85pt}}
\newcommand{\ActiveDigit}[4]{\node[inner sep=0pt] at (#2,#3) {\includegraphics[width=#4mm,height=#4mm,keepaspectratio]{figures/03-foundations/01-basics/tikz/active-learning/assets/digit#1.png}};}
% Both pool views use the same physical canvas and identical sample positions.
% Only the decision boundary and the next queried sample change after retraining.
\newcommand{\ActivePool}[4]{\begin{scope}[shift={(#1,#2)}]
 \node[v3 box,draw=#4,fill=white,minimum width=50mm,minimum height=36mm] (#3) at (0,0) {};
 \begin{scope}
 \clip[rounded corners=1.1mm] (-24.8,-17.8) rectangle (24.8,17.8);
 \ifnum#1=136
  \fill[FigureBlueFill] (-12,-18)--(0,-18)--(-17,18)--(-24,18)--cycle;
  \draw[draw=FigureBlue,line width=.95pt] (-20,18)--(0,-18);
 \else
  \fill[FigureBlueFill] (-18,-18)--(-6,-18)--(-23,18)--(-25,18)--cycle;
  \draw[draw=FigureBlue,line width=.95pt] (-23,18)--(-6,-18);
  \draw[v3 guide] (-20,18)--(0,-18);
 \fi
 \foreach \x/\y/\d in {-20/13/3,-14/8/3,-21/1/3,-14/0/3,-19/-6/3,-11/-8/3,-20/-14/3,-14/-13.5/3,-5/11/8,7/13/8,5/4/8,15/8/8,10/-1/8,17/-5/8,6/-9/8,9/-14/8,-8/3/3}{\ActiveDigit{\d}{\x}{\y}{3.6}}
 \ifnum#1=136
  \draw[draw=FigureRed,dashed,line width=.9pt,rounded corners=.5pt] (-10.8,.2) rectangle(-5.2,5.8);
 \else
  \draw[draw=FigureRed,dashed,line width=.9pt,rounded corners=.5pt] (4.2,10.2) rectangle(9.8,15.8);
 \fi
 \end{scope}
\end{scope}}
\begin{tikzpicture}[v3 picture]
\path[use as bounding box] (0,0) rectangle (163,89);
\node[v3 title] at (24,83) {带标签训练样本};
\node[v3 data,fill=FigureGreen!7!white,draw=FigureGreen,minimum width=44mm,minimum height=21mm] (train) at (24,67) {};
\foreach \x in {8,17,26,35}{\ActiveDigit{3}{\x}{72}{5.2}\ActiveDigit{8}{\x}{62}{5.2}}
\node[v3 note] at (42,68) {$\cdots$};
\node[v3 model,fill=FigureGreenFill,draw=FigureGreen,minimum width=32mm,minimum height=20mm] (net) at (82,73) {};
{\colorlet{VNetworkInput}{FigureGreen}\colorlet{VNetworkHidden}{FigureGreen}\colorlet{VNetworkOutput}{FigureGreen}\VLayerNetTall{82}{73}{1.1}{2}{3}{2}}
\node[v3 title] at (82,86) {重训};
\node[v3 title] at (136,66) {未标注样本池：本轮查询};
\ActivePool{136}{44}{pool}{FigureRed}
\node[v3 title] at (82,28) {人工标注};
\node[v3 output,fill=white,minimum width=39mm,minimum height=21mm] (label) at (82,13) {};
\VMonitor{90}{14.5}{1.4}{FigureBlue}
\node[inner sep=0pt] at (72.5,12) {\includegraphics[width=18mm]{figures/03-foundations/01-basics/tikz/active-learning/assets/annotator.png}};
\draw[FigureMuted,line width=.9pt,line cap=round] (76,6.8)--(101,6.8);
\ActiveDigit{3}{86}{15}{6}\node[v3 label] at (94,15) {$3$};
\draw[v3 flow] (train.east)--(55,67)--(55,73)--(net.west);
\draw[v3 flow] (net.east)--(108,73)--(108,44)--(pool.west);
\draw[v3 flow] (pool.south)--(136,13)--(label.east);
\draw[v3 flow] (label.west)--(58,13)--(58,49)--(24,49)--(train.south);
\node[v3 note,anchor=west] at (49,54) {加入训练样本};
\node[v3 title] at (26,46) {下一轮：边界更新后查询};
\ActivePool{26}{24}{next}{FigurePurple}
\node[v3 note] at (26.5,1.8) {实线：更新后边界　虚线：更新前边界};
\end{tikzpicture}
